\section{Passes}
\label{sec:passes}\label{sec:planning}\label{sec:pass-contracts}
Each subsection gives the pass's signature, its construction rules on the
checked domain, its independent check, and the refinement it supports. Every
constructor checks its target contract. A check $\mathsf{check}_f(s,t,\chi)$
requires both local contracts, compares the ASTs the code compares, and never
calls the constructor. Checks are sufficient structural conditions under the
premises of Section~\ref{sec:premises}: they reject some equal plans and decide
no equivalence.

\subsection{Normalization: CQ $\to$ CQ$_{\rm NF}$}
For an atom at clause $i$, walk its arguments left to right against the declared
types $\tau_j$. $B$ holds names bound by earlier clauses, $L$ names already seen in
this occurrence (reset per occurrence), $A$ bare names this atom supplies that
are not in $B$, and $Q$ is the complete logical scope (including the current
branch's binders). Literals remain in the atom and bind no name. For other
arguments, fresh $k$ starts at \code{__cq_i_j}; cases are tried in order:
\begin{mathpar}
\Rule{A-Variable}{a_j=x\quad x\notin L}
 {a_j\leadsto_{NF}x\quad[L:=L\cup\{x\}]}

\Rule{A-Repeat}{a_j=x\in L\quad k\text{ fresh}}
 {a_j\leadsto_{NF}k\quad[\mathit{after}:=\mathit{after}\concat[k=x]]}

\Rule{A-Wildcard}{a_j=\_\quad k\text{ fresh}}
 {a_j\leadsto_{NF}k\quad[\text{unconstrained witness}]}

\Rule{A-Closed}{a_j=e\quad\fv(e)\cap Q=\emptyset\quad k\text{ fresh}}
 {a_j\leadsto_{NF}k\quad[\mathit{after}:=\mathit{after}\concat[\mathsf{where}\ k==\mathsf{Closed}(e)]]}

\Rule{A-Prefix}{a_j=e\text{ is a remaining expression}\quad
 \fv(e)\cap Q\ne\emptyset\quad\uses(e,A)=\emptyset\quad k\text{ fresh}}
 {a_j\leadsto_{NF}k\quad[\mathit{before}:=\mathit{before}\concat[\mathsf{let}\ k:\tau_j=e]]}

\Rule{A-Local}{a_j=e\text{ is a remaining expression}\quad
 \uses(e,A)\ne\emptyset\quad k,v\text{ fresh}}
 {a_j\leadsto_{NF}k\quad[\mathit{after}:=\mathit{after}\concat[\mathsf{let}\ v:\tau_j=e,\ k=v]]}
\end{mathpar}
Replace the atom by $\mathit{before}\concat[R(\bar a')]\concat\mathit{after}$.
Identities are relational clauses, not guards. Constraints, Rust guards and
lets and generators are copied; branch bodies normalize recursively with fresh names; a
computed head becomes a final typed bind. In \code{R(x,f(x),g())}, closed
\code{g()} remains a once-bound host operand of an occurrence restriction;
the dependent \code{f(x)} bind follows the atom. No closed atom operand creates
a row binding or a Rust block dependency.
\emph{Check:} the NF contract only; there is no separate correspondence check.
\emph{Refinement:} each rewrite replaces an argument by an equivalent
identity or a bind at the position Rust would evaluate it, so meanings are equal.

\subsection{Ordering: CQ$_{\rm NF}\to$ VO}
\label{sec:planning-order}
Let $T$ be scalar-let- and generator-supplied names and $B_w$ the written prefix. Clause $i$ provides
$\bar O_i$ and requires $R_i$:
\begin{mathpar}
\Rule{D-Atom}{c_i=A(\bar x)}
 {B_w;T\vdash c_i\triangleright(\bar x,\set(\bar x)\cap T)}

\Rule{D-Guard}{R=\uses(e,B_w)}
 {B_w;T\vdash\mathsf{if}\ e\triangleright([],R)}

\Rule{D-Bind}{\bar N=\binds(p)\quad R=\uses(e,B_w)}
 {B_w;T\vdash\mathsf{let}\ p[:\tau]=e\triangleright(\bar N,R)}

\Rule{D-For}{\bar N=\binds(p)\quad R=\uses(e,B_w)}
 {B_w;T\vdash\mathsf{for}\ p\ \mathsf{in}\ e\triangleright(\bar N,R)}

\Rule{C-Ready}{i\in U\quad R_i\subseteq B}
 {U;B\vdash i:\mathsf{ready}}
\end{mathpar}
Generators require their RHS free variables, just as scalar lets and guards
do. They are not atom occurrences or relational predicate facts. Later atoms
mentioning their outputs require those bindings rather than supplying them.
There is no written-order chain or predecessor-set condition. Generators
cannot be reduction targets, reduction supports or independent build targets;
a block containing one is not Rust-free. This remains true for a closed
expression or a generator yielding no names.
Identities and constraints require their variables and provide none
(D-Constraint). A Rust-free block provides its free/interface heading; a Rust
block requires its free names and is Rust. Heuristics repeatedly place the
least-ranked ready clause; an empty ready set is an error. With $\kappa$ the
equality-class map and $s_i=|\kappa(O_i)\cap\kappa(B)|$:
\begin{mathpar}
\Rule{S-Written}{}{\mathsf{sequence}_{Written}(q)=[0,\ldots,m-1]}

\Rule{S-Connected}{c_i\text{ atom}\quad s_i=|\kappa(O_i)\cap\kappa(B)|}
 {\mathsf{rank}_{Connected}(i,B)=(1,[s_i=0],0,i)}

\Rule{S-Filtered}{c_i\text{ atom}\quad s_i=|\kappa(O_i)\cap\kappa(B)|}
 {\mathsf{rank}_{Filtered}(i,B)=(1,[s_i=0],f_i,i)}

\Rule{S-Greedy}{c_i\text{ atom}\quad s_i=|\kappa(O_i)\cap\kappa(B)|\\
 r\in\{GreedyWidth,Auto\}}
 {\mathsf{rank}_{r}(i,B)=(1,[s_i=0],|\kappa(B)\cup\kappa(O_i)|,i)}

\Rule{S-Elimination}{c_i\text{ atom}\quad E=\mathsf{analysis}(q).\mathit{elimination\_order}\\
 e_i=\min(\{\operatorname{pos}_E(x)\mid x\in O_i\}\cup\{|E|\})}
 {\mathsf{rank}_{Elimination}(i,B)=(1,0,e_i,i)}

\Rule{O-Coarse}{\mathit{seq}\text{ chosen}}
 {\mathsf{coarse}(\mathit{seq})=[\mathsf{fullSubatom}(i)\text{ or the copied central/Rust step}\mid i\in\mathit{seq}]}

\Rule{O-Fine}{(G_0,B_0)=\mathsf{flush}(B)\quad v\text{ first eligible remaining variable}\\
 (a,G_1,B_1,J)=\mathsf{introduce}(v,B_0)}
 {\mathsf{fineStep}(v,B)=G_0\concat a\concat G_1
   \concat[\mathsf{subatom}(j,B_1)]_{j\in J}}

\Rule{O-BlockPlacement}{b\text{ a block in ordinary ordering}}
 {b.placement=\begin{cases}\mathsf{correlated}&\mathsf{Rust}(b)\\
 \mathsf{hoisted}&\text{otherwise}\end{cases}}

\Rule{O-WrittenBlocks}{t=\mathsf{WrittenCoarse}(s)}
 {\mathsf{compileWritten}_{order}(s)=
 t[\forall b\text{ recursively}.\ b.placement:=\mathsf{correlated}]}

\Rule{O-StepAdvice}{n=N(q)\quad t=\mathsf{authorSteps}(n,A)\quad
 \mathsf{erase}(t)=n\\
 \mathsf{WF}_{VO}(t)\quad
 (L\ne\varnothing\Rightarrow\mathsf{CK\text{-}Le}(n,t,L))}
 {\mathsf{orderAdvice}(q,A)=t}
\end{mathpar}
Ready constraints rank $(0,0,0,i)$, other nonatoms $(0,0,1,i)$; Auto is
GreedyWidth; Elimination uses reverse maximum-cardinality search on the primal
graph. Written/Fine splits atoms in place. $\mathsf{flush}$ emits ready waiting
steps (constraints first); a variable is eligible if no Rust step supplies it or
it is bound; $\mathsf{introduce}$ emits a subatom of its first holding atom and
then the others' checks. Every atom's full correlated row must eventually be
checked. Literal-only occurrences have an empty variable heading and
are placed as complete zero-column reads before the fine variable loop.
The production facade uses Prefix(Connected)/Coarse: the written root, then
at each step a ready atom sharing a bound equality class before a disconnected
one (earliest clause first), with ready relational predicates placed where
their variables are bound. Written/Coarse with O-WrittenBlocks is the exact
executable reference; ordinary Written/Coarse retains O-BlockPlacement and is
the `plan(order(..))` fallback. Filtered/Coarse serves external planners
and R-Masks' sibling order. Under H-FilterRank, $f_i$ is
$0$ for an occurrence containing a literal or a local constant equality, $1$ for a local constant range,
$2$ for a derived domain, relational predicate or finite relational block,
and $3$ otherwise. Equality classes determine locality; Rust guards never
improve a rank, regardless of their syntax or types. This is a syntax-only priority, not a cardinality
estimate. Written remains available to explicit planner clients.

Step advice extends the occurrence selectors with \code{R@i(cols)}: the listed
query variables are the exact retained columns, including every already-bound
column of that occurrence. A bare \code{R@i} reads the full correlated row.
Partial reads introducing a variable are Fine steps (at most one new variable);
adjacent reads of different occurrences on that variable expose intersection
points, without choosing a Physical implementation. A non-introducing explicit
partial read before any Fine read of that occurrence selects it into $L$, the checked
lookup/expand set. Each selected occurrence must then have exactly one later
full expansion with written ownership. An unbound listed variable denotes
a Fine introduction, not an unchecked lookup. Duplicate/foreign columns,
reads after expansion, incomplete rows and unsatisfied dependencies reject.
Ready nonatoms are placed by free-variable readiness, constraints first;
unnamed occurrences follow in written priority, while named partial occurrences
must be completed explicitly. Without step advice, ordinary coarse advice
and the default planner remain unchanged.
\[
\Rule{H-FilterRank}{P_i=\text{local explicit relational predicates on }\kappa(O_i)\quad
 L_i=\{0\mid O_i\text{ has a literal argument}\}}
 {f_i=\min(\{3\}\cup L_i\cup\{\mathsf{strength}(p)\mid p\in P_i\})}
\]
\[
\Rule{CK-Order}{\mathsf{WF}_{NF}(s)\quad\mathsf{WF}_{VO}(t)\\
 \mathsf{OrderObligations}(s,t,\chi)}
 {\mathsf{check}_{order}(s,t,\chi)\Downarrow\ok}
\]
$\mathsf{OrderObligations}$: unchanged CQ; no reductions; every atom has a
subatom; Coarse has one full subatom per atom; Fine introduces at most one
variable per subatom; Written keeps NF order. Filtered/Coarse independently
checks that no ready occurrence outranks the recorded choice and discharges
CK-OrderBranches for the recorded groups. Other rankings remain construction
choices. \emph{Refinement:} by the VO contract (free-variable dependencies,
unchanged total Rust expressions, relational predicates and full rows checked).

\subsection{Hypergraph analysis and GYO}
\label{sec:graph-extraction}\label{sec:gyo-construction}
Edges are atom occurrences and Rust-free blocks over their headings, expanded by
relational equality classes; opaque Rust adds none.
\begin{mathpar}
\Rule{H-Occurrence}{q.body_i=R(\bar x)\text{ or a Rust-free block with heading }\bar x}
 {\mathsf{clauses}(q)[i]=\bar x}

\Rule{H-Incidence}{i\in O\quad x\in\widehat X_i}
 {i\in\mathsf{incidence}(q)[x]}

\Rule{H-Adjacency}{i<j\quad S_{ij}=\widehat X_i\cap\widehat X_j\ne\emptyset}
 {i\xleftrightarrow{S_{ij}}j}

\Rule{G-Private}{i\in\dom Z\quad x\in Z(i)\quad\forall j\in\dom Z\setminus\{i\}.\ x\notin Z(j)}
 {(Z,\pi)\longrightarrow(Z[i\mapsto Z(i)\setminus\{x\}],\pi)}

\Rule{G-Absorb}{i,j\in\dom Z\quad i\ne j\quad Z(i)\subseteq Z(j)}
 {(Z,\pi)\longrightarrow(Z\setminus\{i\},\pi[i\mapsto j])}

\Rule{G-Accept}{Z\text{ terminal}\quad |\dom Z|\le1\quad\forall i\in\dom Z.\ Z(i)=\emptyset}
 {(Z,\pi)\Downarrow\mathsf{acyclic}=true}

\Rule{G-Reject}{Z\text{ terminal}\quad\neg(|\dom Z|\le1\land\forall i\in\dom Z.\ Z(i)=\emptyset)}
 {(Z,\pi)\Downarrow\mathsf{acyclic}=false}

\Rule{H-JoinTree}{\mathsf{acyclic}=true}
 {\mathsf{join\_tree}=\mathsf{Some}([(i,\pi(i))]_{i\in O\text{ increasing}})}
\end{mathpar}
GYO starts at $Z(i)=\widehat X_i$, applies the least private pair, else the least
absorb pair, and records each \code{GyoStep}.
\begin{theorem}[GYO]
GYO terminates and accepts exactly the occurrence hypergraphs with a join tree;
its parent map has running intersection on the original edges.
\end{theorem}
\begin{proof}$|\dom Z|+\sum_i|Z(i)|$ decreases. A variable in two active edges is
never private, so an absorption's shared variables lie in its container: an ear
removal. Both steps preserve join-tree existence, and a join-tree leaf is always
removable, so an acyclic input ends in one empty edge; reversing absorptions
rebuilds the tree.\end{proof}

\subsection{Reduction: VO $\to$ VO}
\label{sec:yannakakis}
Reduction inserts semijoin steps only; it never moves an opaque guard.
With $B=[o,b_1,\ldots]$ the BFS order of the join tree (children by clause index)
and $\mathsf{red}(i,j)$ the reduction of $i$ by $j$ on their shared classes:
\begin{mathpar}
\Rule{R-Guard}{\mathsf{step}=\mathsf{Guard}(e)}
 {\mathsf{reduce}(\mathsf{step})=\mathsf{step}}

\Rule{Y-Up}{B=\mathsf{BFS}(J)}
 {\mathsf{up}(J)=[\mathsf{red}(\pi(b),b)\mid b\in\mathsf{reverse}(B[1..])]}

\Rule{Y-Down}{B=\mathsf{BFS}(J)}
 {\mathsf{down}(J)=[\mathsf{red}(b,\pi(b))\mid b\in B[1..]]}

\Rule{R-Tree}{J=\mathsf{analysis}(q).\mathsf{join\_tree}\ne\mathsf{None}}
 {\mathsf{Tree}(VO)=\mathsf{insertAtReadyPrefix}(VO,\mathsf{up}(J)\concat\mathsf{down}(J))}
\end{mathpar}
The insertion point is the first atom/block, or the earliest prefix providing
the Rust bindings a finite block domain needs. \code{Hypergraph} uses the
filtered reducer only on acyclic scopes, but exact atom-to-block cuts on any
scope: each atom sharing columns with a finite Rust-free block is semijoined
through that block's derived relation. Each branch is analyzed independently.
For R-Filtered, remove a tree edge $(i,j)$ and let $C_{j\mid i}$ be the
component containing support $j$. A candidate $\mathsf{red}(i,j)$ requires
$\min_{k\in C_{j\mid i}} f_k<3$. An upward support first receives its selected
child reductions; siblings run in ascending subtree rank, then written order.
Downward support uses the parent's local filter, sibling subtrees and outside
rank, never a filter learned solely from the target subtree. Explicit
\code{Tree} retains the full reducer for teaching and comparison.
\[
\Rule{R-Filtered}{J\text{ acyclic}\quad
 f_{j\mid i}=\min_{k\in\mathsf{component}(J-\{(i,j)\},j)}f_k\\
 f_{j\mid i}<3\quad\neg(f_{j\mid i}=1\land f_{i\mid j}=0)}
 {\mathsf{red}(i,j)\in\mathsf{filtered}(J)}
\]
Range-only support is omitted against a point-filtered driving component:
its eventual natural join streams the range side instead of prebuilding its
full support index. This optional omission preserves exact joins, makes no
key assumptions, and does not suppress derived-domain support.
Reductions distribute into positive uses and never into negated originals.
\[
\Rule{CK-Reduce}{\mathsf{WF}_{VO}(s)\quad\mathsf{WF}_{VO}(t)\\
 \mathsf{ReduceObligations}(s,t,\chi)}
 {\mathsf{check}_{reduce}(s,t,\chi)\Downarrow\ok}
\]
$\mathsf{ReduceObligations}$: reduction-free source, unchanged CQ, equal step
erasure; under None no reductions; under Tree an acyclic source with every
directed tree edge once, all upward before downward, respecting parent/child
order; Hypergraph requires exact atom-to-block cuts in both cyclic and acyclic
scopes, and only the independently filter-supported tree edges in acyclic
scopes, with strongest upward siblings first. The checker traverses each
cut component independently and checks actual support-version provenance.
Cyclic block-free scopes have no reductions. Missing cuts, unsupported
reductions, and non-tree atom-to-atom reductions are rejected.

The default plan is left-deep in written order: the left natural-join operand
streams and the right builds, and no pass commutes operands. Its reductions
(Masks) follow a reduction tree. H-SpanningCore defines it on every scope: the
GYO join tree when acyclic; otherwise the GYO ears plus a spanning tree of the
residual cyclic core whose edges join clauses sharing a variable (most shared
first, then earliest clauses). Semijoins along it are sound reductions; only on
acyclic scopes are they a full reducer. R-Masks roots that tree at the scope's
first scheduled clause, the outer stream, and reduces each inner clause by its
child subtrees that carry a relational filter or a block cut, children first and
strongest siblings first, before the stream starts; exact atom-to-block cuts
reduce inner atoms. The outer is never a target or a support: its probes filter
it in written order, so no reduced stream is stored. This is the bottom-up pass
of the Yannakakis reducer rooted at the stream, restricted to filter-carrying
subtrees: an unfiltered child subtree has no syntactic evidence that it removes
rows (a key/foreign-key join), so its semijoin is omitted as likely pure
overhead. Consequently Masks is a full reducer only when every edge it omits is
in fact non-reducing on the instance; otherwise a stream row may match an inner
clause and then fail a later unfiltered probe, costing that probe but never
correctness. \code{Tree} remains the unconditional full reducer on acyclic
scopes. CK-Masks checks scope-local edge membership and completeness, not
equality to a reconstructed schedule. Let $C$ be every atom-to-block cut
whose target and support both differ from the outer, and $M$ every
R-Masks edge with an atom target. A subtree is filtered when it contains a
local filter or a target of a cut in $C$. Each edge of $C\cup M$ occurs once
and no other edge occurs. Cuts precede tree edges; every reduction of a
support precedes its use; tree siblings run by subtree filter rank then
written clause. Independent subtrees may interleave. Reductions precede the
first stream step, or, for finite domains requiring Rust-produced columns,
the first source boundary with all those columns available (VO checks
readiness). The outer exclusion includes a block used as the outer.
\begin{mathpar}
\Rule{H-SpanningCore}{T_0=\mathsf{GYOEars}(q)\quad
 T_1=\mathsf{spanningTree}(\mathsf{core}(q),\ \text{most shared, earliest})}
 {\mathsf{reductionTree}(q)=T_0\cup T_1}

\Rule{R-Placement}{E=\mathsf{occurrences}(q)\quad
 C=\{i\mid s_i=\mathsf{block}_{\mathsf{correlated}}(\ldots)\}}
 {\mathsf{reducerGraph}(q,s)=\mathsf{hypergraph}(q,E\setminus C)}

\Rule{R-Masks}{r=\mathsf{outer}(s)\quad c\ne r\quad d\ne r\quad
 d\in\mathsf{children}_{\mathsf{reductionTree}\ \text{rooted at}\ r}(c)\quad
 \mathsf{filteredSubtree}(d)}
 {\mathsf{red}(c,d)\ \text{before the ready stream boundary, children first}}

\Rule{CK-Masks}{\mathsf{edges}(t)=C\cup M\quad
 \mathsf{Once}(t)\quad\mathsf{OuterExcluded}(t)\\
 \mathsf{ReadyBoundary}(s,t)\quad\mathsf{CutsFirst}(t)\quad
 \mathsf{ChildrenFirst}(t)\quad\mathsf{StrongestSiblings}(t)}
 {\mathsf{check}_{reduce}^{Masks}(s,t)\Downarrow\ok}
\end{mathpar}
Correlated blocks denote a prefix-dependent computation, not an independently
available reduction operand. R-Placement excludes their occurrences before
incidence, GYO and tree analysis; their written binders still scope later
clauses. Both reduction construction and checking consume this scope analysis.


\begin{theorem}[Full reducer]
On a join tree, upward reverse-BFS then downward BFS semijoins preserve
$Q=\bowtie_iL_i$ and leave each $L_i=\pi_{X_i}Q$, in $2(n-1)$ steps.
\end{theorem}
\begin{proof}Running intersection makes $X_c\cap X_{\pi(c)}$ the cut of $c$'s
subtree. After the upward pass each node holds the projection of its subtree's
join; the root is fully supported; downward, a surviving child tuple matches a
supported parent tuple and the two witnesses combine. Every full witness
survives every semijoin.\end{proof}
\emph{Refinement:} semijoins are relational and only remove tuples absent from
the join, so meanings are equal; Rust steps are untouched.

\subsection{Elaboration: VO $\to$ RA prime}
State $(D,M,p,B)$: emitted definitions, latest result per clause, optional
prefix, its heading. Relations are emitted lazily at first mention with their
fixed occurrence paths.
\begin{mathpar}
\Rule{E-Emit}{r=\text{next unused }\code{rN}\text{ in emission order}}
 {\emit(D,o)=(D\concat[r=o],r)}

\Rule{E-Relation}{i\notin\dom M\quad C_i=\{c_j=a_j\mid a_j\text{ is literal}\}\\
 m_i:\operatorname{cols}(R_i)\setminus\dom C_i\cong\bar X_i\quad
 (D',r)=\emit(D,\delta_{m_i}(\pi_{\dom m_i}\sigma_{C_i}(R_i@i)))}
 {\mathsf{relation}(i;D,M)=(D',M[i\mapsto r],r)}

\Rule{E-Check}{K\subsetneq X_i\quad\mathsf{introduces}=\emptyset\quad p\ne\bot}
 {(D,M,p,B)\xrightarrow{i,K}(D',M,s,B)\quad[(D',s)=\emit(D,p\ltimes r)]}

\Rule{E-Project}{K\subsetneq X_i\quad\neg(\mathsf{introduces}=\emptyset\land p\ne\bot)\\
 (D',s)=\emit(D,\pi_Kr)}
 {\mathsf{subatomRight}(i,K)=(D',s)}

\Rule{B-Root}{p=\bot\quad s=\mathsf{subatomRight}(i,K)}
 {(D,M,p,B)\xrightarrow{i,K}(D,M,s,B\cup K)}

\Rule{B-Join}{p\ne\bot\quad\text{E-Check inapplicable}\quad s=\mathsf{subatomRight}(i,K)\\
 (D',t)=\emit(D,p\bowtie s)}
 {(D,M,p,B)\xrightarrow{i,K}(D',M,t,B\cup K)}

\Rule{C-Guard}{\mathsf{step}=\mathsf{Guard}(e)\quad p_0=\mathsf{unitize}(p)\\
 (D',s)=\emit(D,\operatorname{guard}_e p_0)}
 {(D,M,p,B)\longrightarrow(D',M,s,B)}

\Rule{C-Bind}{\mathsf{step}=\mathsf{Let}(pat[:\tau]=e)\quad\bar N=\binds(pat)\\
 p_0=\mathsf{unitize}(p)\quad(D',s)=\emit(D,\operatorname{bind}_{pat[:\tau],e}p_0)}
 {(D,M,p,B)\longrightarrow(D',M,s,B\cup\set(\bar N))}

\Rule{C-For}{\mathsf{step}=\mathsf{For}(pat\ \mathsf{in}\ e)\quad\bar N=\binds(pat)\\
 p_0=\mathsf{unitize}(p)\quad(D',s)=\emit(D,\operatorname{for}_{pat,e}p_0)}
 {(D,M,p,B)\longrightarrow(D',M,s,B\cup\set(\bar N))}

\Rule{E-ReduceWith}{\mathsf{step}=\mathsf{reduce}\ i\ \mathsf{with}\ j\\
 (D',s)=\emit(D,M(i)\ltimes_{\equiv}M(j))}
 {(D,M,p,B)\longrightarrow(D',M[i\mapsto s],p,B)}

\Rule{C-Finish}{\text{all VO steps consumed}\\
 p_0=\mathsf{unitize}(p)\quad(D',r)=\emit(D,\pi_Wp_0)}
 {(D,M,p,B)\Downarrow(D',\mathsf{output}\ r\ \mathsf{as}\ h(\bar e))}

\Rule{E-Constraint}{\mathsf{step}=\mathsf{Constraint}(P)\quad
 \vars(P)\subseteq B\quad p_0=\mathsf{unitize}(p)}
 {\mathsf{elaborate}(P,p)=\sigma_Pp_0\text{ at its recorded VO placement}}

\Rule{E-Domain}{C=\bigvee_i\bigwedge_j C_{ij}\quad
 D_{ix}=\{C_{ij}\mid C_{ij}\text{ restricts only }x\text{ against literals or CQ-Closed operands}\}\\
 \forall i.\ D_{ix}\ne\varnothing}
 {C\Longrightarrow D_x=\bigvee_i\bigwedge D_{ix}\quad
  \text{record }D_x\text{ at planning time; retain the exact }C}

\Rule{CT-Constraint}{\vars(C)=\bar x\quad\mathsf{columns}(\bar x)=\bar t\\
 \mathsf{closure}(C)=\{C\}\cup\{D_x\mid\text{E-Domain derives }D_x\}}
 {\mathsf{facts}\longmapsto\mathsf{facts}\cup
   \{\mathsf{predicate}_{shape(D)}(\mathsf{columns}(\vars(D)))\mid D\in\mathsf{closure}(C)\}}

\Rule{RA-Constraint}{\vars(\theta)\subseteq V\quad
 \theta\text{ central and Eq/Ord-admissible}}
 {\sigma_\theta:\operatorname{Set}(V)\longrightarrow\operatorname{Set}(V)}

\Rule{E-BlockPlacement}{b.placement=\mathsf{hoisted}\quad H=F_b\cup I_b
 \quad\text{or}\quad b.placement=\mathsf{correlated}\quad H=B\cup F_b\cup I_b}
 {\mathsf{elaborateBlock}(b,p)=
 \mathsf{arms}(b,H,\mathsf{over}=
   [p\mid b.placement=\mathsf{correlated}])}

\Rule{CK-BlockPlacement}{\forall b\in\mathsf{Blocks}(s).\ \exists u\in\mathsf{Unions}(t).\
 u.origin.path=b.path\\
 \forall b,u\text{ with }u.origin.path=b.path.\
 (u.origin.over\ne\bot)\Longleftrightarrow(b.placement=\mathsf{correlated})}
 {\mathsf{BlockPlacements}(s,t)\Downarrow\ok}
\end{mathpar}
A full subatom always joins; a partial check uses the whole clause relation.
$\mathsf{unitize}$ emits Unit when no prefix exists; binds copy pattern,
annotation and expression. In E-ReduceWith, $\ltimes_\equiv$ matches every
support column with the target's class; when protected Rust or head names block
renaming it is realized as join/select/project. Elaboration consumes recorded
constraint placement (E-Constraint); it does not independently push predicates
or choose joins. Alias representation must preserve logical binding identities.
Recorded hoisted blocks form a derived relation over $F\cup I$, using a named
finite prefix domain for missing generators; recorded correlated blocks run
per row over the prefix. Elaboration never chooses this alternative. The
written-reference ordering entry \code{cq_order::compile_written} records
correlated placement recursively; \code{cq_elaborate::compile_written}
consumes those choices and the written predicate positions.
Structural block views replace every recorded origin of a prefix variable by
the same context column, including equal origins contributed by different
joined atoms. Replacing only the heading representative would leave references
to prefix atoms outside the view.
Query-closed operands are fixed for the invocation just like literal operands.
E-Domain copies their unchanged expression payloads; it does not evaluate or
compare them during compilation. Operand representation therefore cannot remove
necessary per-column domains or their admissible reductions.
The E-Domain closure is recomputed independently while checking both schedules;
unjustified added selections and removal of the original disjunction are
rejected. Explicit relational predicates retain exact constraint and binding
payloads; an opaque source guard cannot be replaced by one because unchanged
Rust expression identity is required. Domain derivation uses only equalities,
comparisons, inclusive ranges and membership over literal or CQ-Closed operands,
never database statistics or foreign-key assumptions.
\[
\Rule{CK-Elaborate}{\mathsf{WF}_{VO}(s)\quad\mathsf{WF}_{RA}(t)\\
 t.declarations=s.cq.declarations\quad t.head=s.cq.head\\
 \mathsf{FixedFacts}(s,t)\quad\mathsf{RustPayloadsAndBindings}(s,t)\quad
 \mathsf{BlockPlacements}(s,t)\Downarrow\ok\quad
 \mathsf{RecordedPlacement}(s,t)\\
 \mathsf{OutputBindings}(s)=\mathsf{OutputBindings}(t)}
 {\mathsf{check}_{elaborate}(s,t,())\Downarrow\ok}
\]
The checker builds a schedule from the supplied VO and compares fixed occurrence
paths and column bindings, equality closure, predicate shapes, named reduction
support, Rust payloads and bindings, output heading and head. Union siblings
match by occurrence footprint, fallback siblings in order, negated references by
original branch. \emph{Refinement:} occurrences, constraints, Rust expressions
and their recorded placements are those of the source; totality and the
representation laws establish equal answers.
For pure positive targets produced by \code{simplify_expansions},
CK-Elaborate also permits an inverse dead-expansion certificate. Starting
only from the supplied RA target, it tentatively replaces a semijoin by a
natural join. Independent demand analysis of that \emph{expanded source}
must prove all added columns unobservable, including newly created implicit
downstream join keys. The ordinary structural comparison then checks the
resulting equivalent witness against the supplied VO schedule. Required
named reductions are still checked on the \emph{original target}; inverse
expansion cannot invent their edges or column matches. This certificate
does not recompile the VO, select a planner, or merge equality facts from
independent existential supports. A split projected-witness view in either
original schedule excludes this fallback: exposing a full read cannot erase
its original correlation obligations. Rust, negative and block programs
retain their existing checker.

\paragraph{Independent dimension branches.}
\textsc{S-JoinBranch} chooses the atom with the most shared equality-class
columns (at least three; written position breaks ties) as a static fact spine.
A locally filtered one-key dimension remains an early independent probe.
The remaining dimension components absorb connected ears only when all
non-spine connections leave through one neighbor heading and the ear's fact-key
set is contained in its parent's. This groups Supplier/Nation and Part/PartSupp without
collapsing a Customer/Supplier cycle. Every member receives the branch's
strongest local filter priority, so a point-filtered Supplier/Nation branch
is scheduled before a merely range-filtered Orders neighbor. These are
compile-time structural heuristics, not estimates of relation cardinalities.
\textsc{WF-JoinBranch} checks the recorded disjoint full-atom groups, their
connectivity and exact external/head interfaces. Opaque Rust, lets and union
blocks conservatively retain the ordinary checked schedule.
\textsc{CK-OrderBranches} checks the widest-spine and pinned-probe premises
independently. Recorded groups must admit connected-ear contractions from
singletons, confined to each recorded group, until each group is one component;
the final partition admits no further ear contraction. Unlisted atoms are
singletons. Every contraction checks shared connectivity, fact-key containment
and containment of all non-pinned external connections in the parent heading.
The checker discharges within-group witnesses by least projected interface
width, earliest parent member, then earliest leaf member. Group record order is
irrelevant. The constructor uses the same priority globally to choose one
deterministic admissible partition, but the checker never invokes that selector
or compares against its output. Other ordering modes forbid recorded groups.
\textsc{E-JoinBranch} elaborates each group to a separate exact natural-join
DAG rooted at its first scheduled atom, with connected child subtrees built
before their parent probes them. It applies local relational predicates,
projects the interface, and joins that named result to the spine. This avoids
a Cartesian product between independent children such as Part and Nation.
No uniqueness or foreign-key premise
is used. \textsc{CK-JoinBranch} independently requires the projected DAG for
each VO group in addition to the unchanged global occurrence, predicate,
reduction and Rust-binding proof. Reduction erasure preserves branch metadata.
For recorded independent branches, the optional Hypergraph policy emits no
semijoin reducer: a branch must not depend on scanning the fact side merely
to build its support. The checker forbids those optional reductions rather
than accepting a missing tree obligation. Explicit Tree on an incompatible
grouped VO is rejected; ungrouped Tree retains the full reducer.
The subsequent RA simplifier retains its existing relational rewrite proof.
\begin{mathpar}
\Rule{S-JoinBranch}{f=\mathsf{WidestSharedKeyAtom}(q)\quad
 B=\mathsf{ConnectedDimensionEars}(q,f)}
 {\mathsf{branches}_{Filtered,Coarse}(q)=B}

\Rule{CK-OrderBranches}{\mathsf{FlatRelational}(q)\quad
 f=\mathsf{WidestSharedKeyAtom}_{\ge3}(q)\quad
 \mathsf{PinnedExcluded}(B,f)\\
 \mathsf{EarProofs}(q,B,f)\quad\mathsf{Saturated}(B)}
 {\mathsf{checkBranches}_{Filtered,Coarse}(q,B)\Downarrow\ok}

\Rule{E-JoinBranch}{b=\{a_1,\ldots,a_k\}\quad I=\mathsf{Interface}(q,b)}
 {\mathsf{elaborateBranch}(b)=
  \pi_I\sigma_{P_b}(R_{a_1}\bowtie\cdots\bowtie R_{a_k})}

\Rule{CK-JoinBranch}{\mathsf{WF}_{VO}(s)\quad
 \mathsf{FixedFacts}(s,t)\quad\mathsf{RustPayloadsAndBindings}(s,t)\\
 \forall b\in s.branches.\ \mathsf{LiveExactProjectedDAG}(t,b,I_b)}
 {\mathsf{checkBranches}(s,t)\Downarrow\ok}
\end{mathpar}


\paragraph{Example: free variables determine eligibility.}
In \code{answer(a,w) :- R(a,x), S(a,y), T(y,z), let k = combine(*x,*y), let
p = left_key(k,*y), let q = right_key(*x,*z), U(p,q,w), where x > 0, if
permitted(k,*z), V(w);} the relational predicate may select $R$ early.
\code{combine} requires only $x,y$, not $T$; \code{permitted} requires $k,z$,
not the preceding independent binding of $q$. Both Rust expressions remain
unchanged. A Rust spelling \code{if *x > 0} would not produce a range read.

\paragraph{Abstract laws.} A ghost join tree witnesses coarse assembly, and
retention justifies projection before a join:
\begin{mathpar}
\Rule{O-Leaf}{i\in O}{i:\{i\}}

\Rule{O-Join}{T_1:C\quad T_2:E\quad C\cap E=\emptyset}
 {(T_1\bowtie T_2):C\cup E}

\Rule{L-Retention}{L:V\quad M:W\quad V\cap(W\cup H\cup\fv(P))\subseteq U\subseteq V}
 {\pi_H\sigma_P(L\bowtie M)=\pi_H\sigma_P(\pi_U L\bowtie M)}
\end{mathpar}
The elaborator projects partial subatoms, block interfaces and domains,
reduction support and the head; it searches no general subtree pruning.

\subsection{Index demands: RA prime $\to$ IX}
\label{sec:pass-index}
Traverse definitions in order from empty $D$:
\begin{mathpar}
\Rule{IX-Reuse}{D(i)=(r,C,supply)}{D\vdash(r,C,supply)\Downarrow(i,D)}

\Rule{IX-New}{(r,C,supply)\notin\operatorname{ran}D\quad i=i_{|D|}}
 {D\vdash(r,C,supply)\Downarrow(i,D[i\mapsto(r,C,supply)])}

\Rule{IX-Join}{o\in\{\bowtie,\ltimes\}\quad
 C=[c\mid 0\leq c<|H(r)|,\ H(r)_c\in H(l)]_{\rm increasing}\\
 D\vdash(r,C,stream)\Downarrow(i,D')}
 {D\vdash j=l\mathbin{o}r\Downarrow D'}

\Rule{IX-Const}{\theta\text{ is a point conjunction on columns }K\\
 \text{each term is a constant equality or literal membership}\\
 C=\operatorname{ordinals}_{H(r)}(K)\quad
 D\vdash(r,C,const)\Downarrow(i,D')}
 {D\vdash j=\sigma_\theta r\Downarrow D'}

\Rule{IX-Range}{\theta\text{ conjunction of constant/parameter order bounds on }c\\
 C=[\operatorname{ordinal}_{H(r)}(c)]\quad
 D\vdash(r,C,range)\Downarrow(i,D')}
 {D\vdash j=\sigma_\theta r\Downarrow D'}

\Rule{IX-Other}{d\in\{\mathsf{Unit,Rename,Project,Guard,Bind,Union}\}\\
 \text{ or a select matching neither IX-Const nor IX-Range}}
 {D\vdash d\Downarrow D}

\Rule{CK-Index}{\mathsf{WF}_{RA}(s)\quad\mathsf{WF}_{IX}(t)\quad t.ra=s\\
 t.indexes=[i_n=d_n]_{n=0}^{k-1}\text{ from first-use derivation}}
 {\mathsf{check}_{index}(s,t,())\Downarrow\ok}
\end{mathpar}
Antijoin uses IX-Join on its explicit key only. \emph{Refinement:} equality,
since IX adds no meaning.

\subsection{Physical lowering: IX $\to$ Physical}
Lowering restates each RA result with one physical operator and the mode its
reads allow; defaults are first probes, eager semis and anti, union defaults of
Section~\ref{sec:physical}.
An occurrence's literal bindings remain on its rename, not a select operator.
Such a rename is retained (\code{let}, or \code{exists} for an empty heading):
every chosen index is built over the filtered, projected occurrence, and a
stream source retains filtered rows rather than filtering a reusable host scan.
\begin{mathpar}
\Rule{PH-Range}{\theta\text{ constant order bounds on one column }c}
 {\mathsf{reads}(range\ select\ r\ if\ (\theta))=(r,\{c\},range)}

\Rule{CK-Demands}{\mathsf{WF}_{Physical}(t)\quad
 \forall(r,C,supply)\in\mathsf{keyedReads}(t)\text{ (excluding merge-left)}.\\
 (r,\operatorname{ordinals}_{H(r)}(C),supply)\in\operatorname{descs}(s)}
 {\mathsf{demands}(s,t)\Downarrow\ok}

\Rule{CK-Physical}{\mathsf{WF}_{IX}(s)\quad\mathsf{WF}_{Physical}(t)\\
 \mathsf{erase}(t)=s.ra\quad\mathsf{demands}(s,t)\quad
 \textsc{CK-Return}(t,\Gamma)}
 {\mathsf{check}_{physical}(s,t)\Downarrow\ok}
\end{mathpar}
No missing demand is synthesized; merge has its own ordering premise; the check
accepts any admissible recorded decisions. \emph{Refinement:} erasure equals the
source RA, and modes only change evaluation.

\subsection{Structure selection: Physical $\to$ Structures}
The default constructor copies Physical and records default types for every
built level, host reader fallback, retained collection and dedup boundary,
including support for the selected keyed layouts;
compatible membership prefixes may share a curried structure; a smallest
intersection records two candidate levels per retained side; const/range
supplies never justify a build. This restriction concerns predicate reads,
not atom literals: literal-bearing renames always build their filtered relation.
Their entries use the ordinary \code{let} retention and index rules, without
filtered-storage annotations. Literal restriction and projection belong only to
the occurrence's row source; GPU materializes that source in the build phase.
The optional \code{positional(H,...)} arguments assert host capabilities at this
boundary; Physical contains only relational schemas, so unspecified hosts keep
the previous representation. The initial selector handles recorded point/range
reads and single-column eager membership; Rust capability dispatch chooses the
key codec without compiler inspection of Rust types. It chooses retained masks
for eligible rescans or shared reads and leaves single filtered index-only
builds streaming. Supplemental exact keysets beside payload indexes are explicit
entries. Multi-key reductions, lazy semi and unsupported spines keep their old
representation. Selection point constructors currently require a single exact
point conjunct; positional pipe reads retain all residual conjuncts.
\begin{mathpar}
\Rule{SEL-Positional}{H\in\mathsf{positionalArguments}\quad
 \mathsf{lineage}(r)=(H,\lambda)\quad\mathsf{eligibleReads}(r)}
 {\mathsf{select}(r)\text{ records base, lineage, marker, provider and summary}}
\end{mathpar}
\[
\Rule{CK-Structures}{\mathsf{WF}_{Physical}(s)\quad\mathsf{WF}_{Structures}(t)\\
 s=t.physical\text{ as AST and serialized tokens}}
 {\mathsf{check}_{structures}(s,t)\Downarrow\ok}
\]
\emph{Refinement:} the embedded Physical is unchanged; each type
represents its relation under the storage laws (a premise checked by tests).

\subsection{Realization: checked Sequence $\to$ Rust}
The default CPU translator consumes checked, sequentially graded Sequence.
Its host interpreter uses row adapters over the \code{cq-storage} traits;
the shared realization helpers also serve an external push printer that
consumes Structures directly. Rust translation has no source/target
correspondence check; evidence is answer oracles, rustc trait guards and
expansion fixtures. Interfaces at the Rust boundary:
\begin{mathpar}
\Rule{IF-Surface}{\operatorname{parse}_{CQ}(T)=s\quad
 \text{no unexpanded extension clause remains in }s}
 {\mathsf{surface}(T)=s\text{ is admitted to CQ checking}}

\Rule{IF-Input}{\text{unique relation and column names}\quad\text{positive arity}}
 {\mathsf{layout}(d)=\operatorname{columns}(d)\text{ in declaration order}}

\Rule{IF-Output}{A=\mathsf{collect\_from}(I)\quad
 \alpha(A)=\{r\mid r\in I\}}
 {A:\code{Canonical}\quad A\text{ may supply an input host}}

\Rule{IF-Environment}{\mathsf{visibleLogicalUses}(e)\subseteq\Gamma_Q\quad
 \Gamma_H,\Gamma_Q\vdash_{Rust}e}
 {e\text{ is embedded Rust in the closing environment}}

\Rule{IF-Parameter}{p\in\vars_H(\theta)\quad b_p\text{ fresh}}
 {p\longmapsto\code{let b_p = \&p;}\text{ once outside all query loops}}

\Rule{IF-Range}{C=\{d_1\}\quad H:\code{OrderedRead}\text{ at the matching key type}}
 {\mathsf{window}(H,C)\text{ is an ordered read; otherwise scan/filter, no build}}

\Rule{IF-Leaf}{d=\code{relation R(...)=}(e)\quad b\text{ fresh}\quad
 \Gamma_H\vdash_{Rust}e}
 {d\longmapsto\code{let b = \&(e);}\quad R\mapsto b}

\Rule{IF-Dispatch}{\mathsf{nativeEligible}(H,C)\quad H:\mathsf{ReadAt}(C)}
 {\mathsf{read}(H,C)\longmapsto\mathsf{nativeRead}(H,C)}

\Rule{IF-Fallback}{\neg\mathsf{nativeEligible}(H,C)\ \lor\ H\not:\mathsf{ReadAt}(C)}
 {\mathsf{read}(H,C)\longmapsto\mathsf{lazyOnce}(\mathsf{build}(\alpha(H),C))}
\end{mathpar}
For keyed host reads, $\mathsf{nativeEligible}(H,C)$ holds at a leading
declaration prefix, or at a single nonleading borrowed-key column
$C=\{d_{j+1}\}$ when $H:\code{IndexedColumn<j>}$. The zero-based slot is derived
from the checked declaration and rename, not inferred from matching Rust types.
The explicit read's key is the one-field tuple of that column's declared owned
type; its payload supplies every other column in declaration order, with exact
source occurrences and chronology. Its cursor borrows or pins that same
immutable source, not a temporary dispatch wrapper. Exact membership requires
its separate Rust trait; retention and expansion still require Copy and
Copy plus ExactSizeIterator respectively. An unsupported slot or capability
uses the ordinary fallback. Computed keys, multicolumn nonleading keys,
candidate enumeration and PointLookup retain their existing paths.

\begin{mathpar}
\Rule{IF-Positional}{r\text{ has a checked positional read }(H,C,P)}
 {r\longmapsto\code{BlockRead::read}(H,C,P)\text{ then survivor row fetch}}

\Rule{IF-Selection}{r\text{ records an immutable selection and its read spine}}
 {r\longmapsto\code{Selection::from_positions/narrow};\quad
  \mathsf{scan}(r)\longmapsto\code{rows};\quad
  \mathsf{index}(r)\longmapsto\code{build}}
\end{mathpar}
The recorded predicates emit \code{Point}, \code{Window} or
\code{Membership}; residual conditions precede further reducers. The additive
\code{BlockRead::read_owned} form owns a pipe predicate and cursor state, so
no local predicate reference forces output collection. It shares the same
bounds and exact read law as borrowed \code{read}; neither form requires cloning
keys or predicates. A pipe keeps a lazy ID iterator; only a recorded retained
selection builds an immutable mask, consuming the complete checked read spine once.
Piped predecessors never require intermediate snapshots. Keysets consume only
recorded keys through Build; their provider performs exact dense/hash adaptation
within its recorded budget. Each named payload index uses the snapshot's Build
projection, preserving constants, multiplicity and original host references.
No tuple buffer is inserted beside a selection. Narrowing a retained predecessor
enumerates its set bits in each surviving block, not every base position, and
uses the named selection's summary policy. Empty blocks and blocks rejected by
enclosing bounds are skipped; partial and non-word-aligned block boundaries
preserve occurrence order. Skipping is in memory after decoding.
Lazy semi keeps the incremental KeySet route and is not drained for summaries;
opaque, branch and unsupported inputs retain their previous demand behavior.
No nonlocal probe-filter motion is implemented.

\paragraph{Push realization.}
The push printer consumes the same checked Structures decisions and shared
source preparation. Preparation exposes the fused positional pipe before any
later iterator adapter consumes its compiler tokens. Push loops over those
accepted IDs, fetches each surviving host row, and invokes its sink directly.
All residual predicates remain inside the prepared ID stream. A retained let
constructs its named Selection once; repeated scans use that snapshot's rows,
and payload indices and AdaptiveKeySets use the recorded Build projections.
There is no intermediate traversal Vec beside a Selection. Key-only membership
(including anti and computed-left reads) uses the recorded ownership layout;
negative membership does not use positive interval rejection. Standalone union
arms keep their scopes, and fused dependent branches consume their prepared
stream once. Lazy membership retains its existing empty-left nonforcing rule.
\begin{mathpar}
\Rule{IF-PushPositional}{r\text{ has checked ID stream }I\text{ over base }H}
 {\mathsf{push}(r,k)\longmapsto
  \code{for id in I \{ k(row\_at(H,id)); \}}}
\end{mathpar}
The user advises \code{select(positional(H,...))}; the selection pass decides
and records the representation. Without advice the existing representation is
retained. Both printers enforce the Structures contract and emitted
\code{PositionalRows}/\code{BlockRead} guards. Neither infers positional
capabilities from a Rust host type. Push continues to reject written merge and
smallest choices, including unused results.

Native keyed reads are chosen by stable autoref dispatch at the exact
leading-prefix key type, with the recorded fallback otherwise; const/range reads
scan and filter without building; merge requires \code{OrderedRead}. Rows borrow
host and built fields; computed columns and lookup keys are owned; canonicality
never follows from a host type.

\paragraph{Host-set certificates.}
\begin{mathpar}
\Rule{IF-HostSet}{\code{host\_set R}\in t\quad
 \mathsf{exactDistinct}(\mathsf{fullRows}(R))=true}
 {\mathsf{certify}(R)\Downarrow\ok}
\end{mathpar}
Before any operator relies on \textsc{PH-HostSet}, the printer checks the
entire input heading, not a query projection. Duplicate rows fail the query
with an error naming the relation; hash collisions are resolved by exact
tuple equality, never by fingerprints alone. CPU printers share eager
full-row certification. The GPU printer emits an exact device certificate:
adjacent comparison simultaneously checks lexicographic sortedness and
duplicates; unsorted input uses an exact row-index hash table. Only a
successful certificate permits query execution. This obligation is absent
without the opt-in fact and leaves the default CPU path unchanged.

\paragraph{Opaque generator realization.}
Write $I_r$ for the input row stream and $\mathsf{extend}_p(\rho,v)$ for
destructuring item $v$ and extending the row with the fresh pattern columns.
The same checked Bind node has the following two realizations:
\begin{mathpar}
\Rule{P-For}{\mathsf{Bind}(r,p,e)\text{ carries the \code{for} token}\\
 \mathsf{scopeOK}\quad\mathsf{unchangedExpression}\quad\mathsf{finitePureTotal}(e)}
 {\mathsf{pull}:
 I_r.\mathsf{flat\_map}(\rho\mapsto
 \mathsf{into\_iter}(e(\rho)).\mathsf{map}(v\mapsto\mathsf{extend}_p(\rho,v)))\\
 \mathsf{push}(k):
 \mathsf{for}\ \rho\ \mathsf{in}\ I_r\
 \{\mathsf{for}\ v\ \mathsf{in}\ e(\rho)\
 \{k(\mathsf{extend}_p(\rho,v));\}\}}
\end{mathpar}
Both evaluate the RHS in the input environment, exactly at the opaque step;
iteration stays inside that step's row scope. Generated item columns have
the same owned-computed representation as scalar bind columns, while existing
host references retain their origins. No collector, sorting or deduplication
is inserted at the generator itself; existing canonical boundaries suffice.
Shared generator callables retain host-column lifetimes for iterators borrowing
their inputs, but quantify computed-column borrows independently for every
call. A row-local computed value must not borrow for the enclosing callable's
lifetime merely because the unchanged expression also captures host state.
Zero-sized type witnesses connect prepared scalar results and generator item
patterns to later callable inputs. Rust infers unannotated computed columns
from these connections; preparing a witness evaluates neither the original
expression nor its pattern-extractor closure.
Absent the token, the scalar printer remains a one-row extension and must not
call \code{IntoIterator}. Generator iteration is pure and total by premise;
stage eligibility depends on its representation. Rust checks iterator/item
and lifetime obligations; the checker preserves the tag, original expression
and free-variable bindings; finite pure total iteration remains a host premise.

\subsection{Shared planning compositions}
With $N$ normalization, $O_{h,s}$ ordering, $R_{g,r}$ reduction, $E$ elaboration,
and $S$ checked redundant-support absorption, Plan-Prefix is the reusable
\code{cq_planning::compile_prefix} core composition, re-exported by
\code{mini_linq::planning} and \code{cq_planners_plan}, not a macro-local policy.
Its \code{PlanOptions} parser is shared by macro and programmatic callers:
\code{order(...)} and \code{no_reduce} occur at most once, reduction and skip
packets may repeat, and a trailing options comma is rejected. The checked
\code{prefix_valuation} API exposes $w$ before elaboration.
Qualified \code{PlanOptions} selectors name clauses of the authored input $q$.
Normalization preserves top-level atom/block order, so its inserted bindings
and restrictions induce a positional map $\rho$ from those source clauses
to $N(q)$. Apply $\rho$ to qualified order, reduction and skip selectors before
passing them to the normal-IR ordering API; unqualified selectors are unchanged.
\begin{mathpar}
\Rule{Plan-Prefix}{p=\rho(\mathsf{parsePlanOptions}(args))\quad n=N(q)\\
 v=\begin{cases}O_{explicit(p.order),Coarse}(n)&p.order\text{ present}\\
 O_{Prefix(Connected),Coarse}(n)&\text{otherwise}\end{cases}\\
 r=\text{None if }p.no\_reduce,\text{ otherwise Masks}\\
 w=\mathsf{advise}(R_{None,r}(v),p)\quad a=E(w)}
 {\code{compile_prefix}(q,p)=\code{mini_linq::passes::plan}(q,args)=a}

\Rule{R-Advice}{\mathsf{skip}(T,S)\Rightarrow\mathsf{red}(T,S)\in w\quad
 \mathsf{reduce}(T,S)\ \text{or}\ \mathsf{skip}(T,S)\Rightarrow
 T\ne S\ \text{top-level atoms},\ \mathsf{shared}(T,S)\ne\emptyset}
 {\mathsf{advise}(w)=w\setminus\{\mathsf{red}(T,S)\mid\mathsf{skip}\}
  \ ;\ \{\mathsf{red}(T,S)\mid\mathsf{reduce}\}\ \text{before the stream}}

\Rule{Plan-Generic}{n=N(q)\quad v=O_{Written,Fine}(n)\quad a=E(v)}
 {\code{cq_planners::order(written, fine)}(q)=a}

\Rule{Plan-Graph}{n=N(q)\quad v=O_{Written,Coarse}(n)\\
 r=\text{Hypergraph if blocks; otherwise Tree if acyclic, else None}\\
 w=R_{None,r}(v)\quad a=S(E(w))}
 {\code{cq_planners::order(hypergraph)}(q)=a}

\Rule{Plan-Transfer}{n=N(q)\quad v=C_{atoms}(O_{Written,Coarse}(n))\quad w=R_{None,Masks}(v)\\
 r=\mathsf{outer}(w)\quad T=\mathsf{reductionTree}(q)\ \text{rooted at}\ r\\
 F=\{\mathsf{red}(r,c)\mid c\in\mathsf{children}(r),\ \mathsf{filteredSubtree}(c)\}\\
 B=\{\mathsf{red}(c,p)\mid p=\mathsf{parent}(c),\ \exists\ \text{filter outside}\ \mathsf{subtree}(c)\}\ \text{parents first}}
 {\code{cq_planners::order(transfer)}(q)=S(E(w;F;B))}

\Rule{CK-Transfer}{\mathsf{erase}(t)=v\quad t=\mathsf{transfer}(R_{None,Masks}(v))}
 {\mathsf{check}^{Transfer}(q,t)\Downarrow\ok}

\Rule{Plan-Le}{n=N(q)\quad v=O_{Written,Coarse}(n)\quad w=R_{None,Masks}(v)\\
 \text{for each inner occurrence }R\text{ in GYO ear order:}\\
 \mathsf{lookup}(R)=R(\bar b,\_)\ \text{at the first point its bound key }\bar b\text{ is available}\\
 \mathsf{expand}(R)=R(\bar b,\bar y)\ \text{at the latest point before }
 \mathsf{firstUse}(\bar y)\lor\text{scope exit}}
 {\code{cq_planners::order(le)}(q)=E(\mathsf{le}(w))}

\Rule{CK-Le}{\mathsf{erase}(t)=n\quad
 \forall R\in L.\ \mathsf{lookup}(R)<\mathsf{expand}(R)\ \text{exactly once each}\\
 \mathsf{columnsOwned}_{L}(t)=\mathsf{columnsOwned}_{L}(n)\quad
 \mathsf{outer}(t)=\mathsf{outer}(w)\\
 \text{WF-VO}(t)}
 {\mathsf{check}^{Le}(q,t,L)\Downarrow\ok}
\end{mathpar}
The \code{no_reduce} option disables implicit Masks reductions, not explicit
\code{reduce(T with S)} advice. Such advice therefore composes with the empty
Masks schedule. A \code{skip(T with S)} still requires an existing Masks
reduction; under \code{no_reduce} it fails the R-Advice premise rather than
silently becoming a no-op.

\code{order(transfer)} is an exact-semijoin adaptation of robust predicate
transfer (Zhao et al., SIGMOD'25), not the paper's vectorized Bloom pipeline or
RPT+'s adaptive execution. $C_{atoms}$ preserves every predicate, binding and
block boundary, but within each maximal top-level run of atoms prefers the
earliest atom connected to already bound columns. It preserves written ties
and does not invent a connection for genuinely disconnected components.
This prevents avoidable Cartesian prefixes even when opaque Rust excludes
additional transfer reductions. The reduction tree is rooted at the stream;
when that is the largest relation the root coincides with LargestRoot.
The most-shared-first spanning core matches the paper's maximum-spanning-tree weights.
The forward pass is the Masks pass plus the stream's reductions by its filtered
children; the backward pass then reduces each child atom by its already reduced
parent. As in RPT, no transfer starts from a side without filters, transfers
follow only edges whose clauses share a variable (a join predicate), and blocks
are never transfer supports beyond the Masks atom-to-block cuts. Transfers are
exact semijoins, not Bloom filters. Checked absorption removes redundant support
versions before Physical chooses which remaining results to retain. Scopes containing
opaque Rust keep only their Masks reductions: there the structural RA
correspondence check rejects some sound transfer schedules.

\code{order(le)} is the Lookup \& Expand decomposition of Birler, Kemper and
Neumann (VLDB'24) as a VO construction: a lookup is a partial subatom that
introduces no names (a semijoin of the stream prefix), an expansion is the
later full subatom of the same occurrence. Lookups are placed in GYO ear order
as soon as their bound key is available, expansions as late as column demand
or scope exit permits, so acyclic scopes form two-phase plans.
Expansions are coarse full rows under an explicit CK-Le obligation, not Fine
subatoms. Whether a later probe reuses the lookup's cursor is a Structures
representation decision, not part of this plan.
The phase checker is shared by \code{cq_order} and the external planner, rather
than reimplemented by the facade. For Plan-Le, $L$ is every inner occurrence
in each scope, the outer has one full read with written ownership, and the
caller additionally checks Masks reductions. For step advice, $L$ is only the
explicit lookup occurrences at the top level; unselected Fine/coarse reads and
unchanged nested scopes retain WF-VO obligations. CK-Le then requires the
written outer stream whenever $L$ is nonempty. The phase check runs before
optional reduction advice.

\code{plan} is left-deep in written order and defaults to checked Masks
reduction: every scope, acyclic or cyclic, reduces its inner clauses along the
outer-rooted reduction tree (R-Masks) and by exact atom-to-block cuts; every
branch is its own scope. Reduction compilation is independently checked before
elaboration, and elaboration is checked against the reduced VO schedule. The
external Hypergraph and Transfer planners additionally apply a checked RA
realization rewrite. It uses
$L\ltimes(R\ltimes L)=L\ltimes R$ and
$L'\bowtie(R\ltimes L)=L'\bowtie R$ only when the required shared columns of $L'$
are proven contained in $L$ by reduction lineage. The corresponding left-operand
law removes $(L\ltimes R)\bowtie R$'s duplicate membership test (and similarly
for a semijoin consumer). Right-side support pruning takes precedence; the
left-side law must be proved using the surviving RHS, not a discarded
restriction. This is local absorption, not a maximal fixed-point guarantee.
Dead reduced copies are removed; independent consumers that still need a
reduction keep the same restricted witness alive.
The rewrite is checked against the original RA and does not alter the reduction
algorithm.

For LE's checked RA \code{simplify_expansions} pass, let $H_L,H_R$ be the
source operand headings, $K=H_L\cap H_R$, and $D_s(r)$ the demanded columns of source result
$r$. Under set canonicalization the additional local laws are
\[
 \pi_D(L\bowtie R)=\pi_D(L\ltimes R)\quad(D\subseteq H_L),
 \qquad L\ltimes R=L\quad
   \bigl(\pi_K L\subseteq\pi_K R\bigr).
\]
Thus a named expansion $r=L\bowtie R$ may become $L\ltimes R$ only when
$D_s(r)\subseteq H_L$; with the same demand premise it may be bypassed when
the support is already witnessed. A redundant semijoin may likewise be
bypassed. These laws concern canonical results, not bag multiplicities or
opaque Rust invocation counts.

The certificate independently propagates demand backwards through the
\emph{source} DAG: seed the output heading, retain projection attributes and
selection variables, and require every natural-join and semijoin shared key
even if its column is absent from the final output. Propagate these demands
to each operand on its source heading. Merely checking that the target is
well headed is insufficient: dropping a column can silently remove a later
implicit join equality. Support containment requires the \emph{same named
result} $R$ in $L$'s positive lineage, with every column of $K$ retained at
each step. Lineage may follow projection, selection, either natural-join
operand, or either semijoin operand: a semijoin guarantees an existential
right witness, but only on keys retained along the complete path.

The certificate replays these local laws, recomputing source demand after
each proved step, and finally requires the exact proposed target. It
substitutes proved aliases in every surviving operand and output and permits
removal only of unreachable pure definitions. The new laws conservatively
decline programs containing binds, guards, anti or union; opaque execution,
negative originals and block ancestry are therefore unchanged. The general
schedule contract is not weakened.
Set
\code{pipeline \{ plan = mini_linq::passes::plan(no_reduce); \}}
to disable optional semijoins while keeping the default Prefix(Connected)
order. Configuration
passes this argument through its existing named planning invocation.
\code{plan} uses the optimizing elaborator, so even the opt-out is not
the exact reference. Generic and the external Graph spelling belong to the
external \code{cq-planners} consumer, which preserves its explicitly requested
ordering policy.

\paragraph{Join-filter selection and realization.}
\textsc{SEL-JoinFilter} records a scalar filter for every admissible
(\textsc{STR-JoinFilter}) covering probe or semijoin into a statically filtered,
total single-level build, one per scalar column of its key, at the scan its
exclusive spine reaches; this is the bitmap a build side pushes to the deepest
probe-side scan. A scan may therefore carry filters from several builds.
Every eligible payload-index build selects an exact
\code{AdaptiveMembership<K,67108864>}. A membership-only build already answers
the same exact key test and receives no redundant supplementary filter.
Rust's \code{DenseOrdinal} capability selects the dense representation under
\textsc{STR-AdaptiveMembership}; other keys keep exact hash membership.
Selection does not sample cardinalities or assume keys or foreign keys.
Approximate structures live in an external feature and require explicit
\code{cq_rewrites::retype} advice; the printer has no algorithm-specific branch.
Every recorded covering witness is validated before runtime filters are shared.
Sharing requires the same scan occurrence, build version, scalar key, source-index
layout, filter type and column marker; covering-operator identity alone does not
require another filter. First-witness order supplies scan-local identities, also
used to distinguish profile labels.
\textsc{IF-JoinFilter} requires recorded \code{keys} support on the selected index
layout and constructs each distinct checked provider from that index's
\code{Keys::distinct_keys}, projected to the filtered key column, through
\code{cq_storage::build!} on $(k,())$ pairs. Rust proves the marker's
\code{ColumnSlot} and \code{require::conservative::<K,T>} obligations, then realization
filters positions before fetching payloads: the first filter through the
block-skipping read, the others on its surviving positions, in decision order.
Host block bounds use the provider's conservative \code{any_in}; an absent bound
never skips a block. Each covering exact operator is unchanged, so conservative false
positives only perform extra probes. No opaque guard or bind is crossed.
A scan read before some covering provider is built (for example a retained
host publication ordered ahead of that build) applies only the filters already
built: an omitted filter is the trivial conservative filter that passes every
key, so the covering exact operators still decide membership.
\textsc{IF-JoinFilterAbandon} decides at runtime whether to construct the
filters of a build at all: a build whose \code{Keys::key_count} exceeds 35\% of the rows
of its positional source occurrence constructs none (the creation threshold of
dynamic predicate transfer, RPT+), and its abandoned filters pass every key,
which the one-sided conservative law permits. A build without a single
positional source occurrence always constructs its filters. Only the covering
exact operators decide membership either way.
\begin{mathpar}
\Rule{SEL-JoinFilter}{
 \mathsf{eligibleCoveredScan}(s,b,k,j)\quad
 \mathsf{filteredBuild}(b)\quad\neg\mathsf{membershipOnly}(b)\quad
 \mathsf{positionalHost}(s)\\
 T=\code{AdaptiveMembership<K,67108864>}
}{
 \mathsf{select}(s,b,k,j)=
 \mathsf{joinFilter}(T,b,k,j,\mathsf{marker}(s,k))
}

\Rule{IF-JoinFilter}{
 F_i=\mathsf{buildFilter}_{T_i}(\pi_{k_i}\mathsf{keys}(\mathsf{index}_{b_i}))\quad
 \mathsf{STR\text{-}JoinFilter}(s,F_i,b_i,k_i,j_i)\quad
 \mathsf{builtBefore}(b_i,s)\ (1\le i\le n)\\
 \mathsf{keepBlock}(B)=
 \mathsf{staticMayMatch}(B)\land F_1.\mathsf{anyIn}(\mathsf{bounds}_{k_1}(B))
}{
 \mathsf{scan}(s)=
 \mathsf{filter}_{\bigwedge_i F_i.\mathsf{test}(k_i)}
 (\mathsf{staticRead}_{\mathsf{keepBlock}}(s))
}

\Rule{IF-JoinFilterAbandon}{
 \mathsf{source}(b)=s'\quad
 20\,|\mathsf{keys}(\mathsf{index}_b)| > 7\,|s'|
}{
 F=\top\quad(\forall x.\ F.\mathsf{test}(x))
}
\end{mathpar}
Absent bounds make the corresponding block test true. The emitted filter and
exact probe preserve surviving occurrence order and multiplicity.
